% ============================================================================
% ICSE Paper — Approach Section (v3b: fused prose, succinct overview)
% LLM-Driven Trace Link Recovery between Architecture Documents and Models
%
% Changes from v3:
%   - Challenge/Insight/Design labels removed; content flows as natural prose
%   - Challenges reorganized: knowledge gap is the parent challenge,
%     reference diversity and name ambiguity are its two manifestations
%   - C4 (non-determinism) and C5 (overlap) dropped as standalone challenges;
%     non-determinism is woven into subsections, overlap into consolidation
%   - Challenge table removed (too fine-grained for 3 challenges)
%   - Overview shortened (~50%); subsections remain detailed
%   - Plain English throughout for non-native speakers
% ============================================================================

\section{Approach}\label{sec:approach}

\subsection{Overview}\label{sec:overview}

We propose an LLM-driven pipeline for recovering trace links between
Software Architecture Documents (SAD) and Software Architecture Models
(SAM).
Given a natural-language architecture document with sentences
$S = \{s_1, \ldots, s_n\}$ and a component model with components
$C = \{c_1, \ldots, c_m\}$, our approach produces a set of trace links
$L \subseteq S \times C$, where each link $(s_i, c_j)$ asserts that
sentence~$s_i$ discusses component~$c_j$.

Two problems make this task difficult.
First, \emph{reference diversity}: documents refer to components by
canonical name, abbreviation, functional description, trailing word,
or pronoun.
No single extraction strategy handles all of these forms.
Second, \emph{name ambiguity}: component names like \texttt{Cache} or
\texttt{Facade} are common English words that appear frequently
without referring to any component, so broadening extraction to
capture diverse references also increases false positives.

Addressing both problems requires structured knowledge---which names
are ambiguous, and what abbreviations, synonyms, and short forms map
to which component.
This knowledge does not exist in advance and must be discovered at
runtime, using the same LLM that will consume it, with the risk of
hallucinated mappings or misjudged ambiguity.
We call this the \emph{knowledge gap}.

Existing approaches address only part of this problem.
Lexical NLP pipelines such as SWATTR~\cite{swattr} rely on string
similarity and cannot resolve indirect references or discourse
dependencies.
LLM-based classifiers such as LISSA~\cite{lissa} add semantic
understanding but have no mechanism to build alias mappings or classify
name ambiguity before linking, treat each sentence in isolation, and
do not address LLM non-determinism.

Our pipeline closes the knowledge gap first, then recovers links
using the acquired knowledge.
It is organized as a DAG (Figure~\ref{fig:pipeline}) with three
layers:
\textbf{Knowledge acquisition} (Layer~1) extracts component name
properties and an alias table from raw inputs
(\S\ref{sec:knowledge});
\textbf{Link recovery} (Layer~2) runs four complementary strategies
targeting different reference types
(Table~\ref{tab:strategies}, \S\ref{sec:recovery});
\textbf{Link consolidation} (Layer~3) deduplicates overlapping
candidates by first-seen priority (\S\ref{sec:merge}).
To handle LLM non-determinism, the pipeline uses two cross-cutting
patterns.
\emph{Multi-pass consensus}: precision-sensitive decisions use two
independent LLM passes with task-specific agreement rules.
\emph{Citation-grounded verification}: the LLM must produce
checkable evidence---such as antecedent sentence numbers---that
deterministic code can verify.
To prevent data leakage, all few-shot examples use abstract
placeholders or safe textbook domains; no benchmark names or
gold-standard information appears in any prompt.

\begin{table}[t]
\caption{Link recovery strategies and their properties.}
\label{tab:strategies}
\centering\footnotesize
\begin{tabular}{@{}lllll@{}}
\toprule
\textbf{Strategy} & \textbf{Reference type} & \textbf{Knowledge} & \textbf{Consensus pattern} & \textbf{LLM} \\
\midrule
Explicit & Canonical name & Alias table\rlap{$^\ddagger$} & Asym.\ voting + stratified valid. & 2+2--3\rlap{$^*$} \\
Contextual & Alias, functional & Alias table, ambig.\ classif. & $\cap$ extr.\ + stratified valid. & 2+2--3\rlap{$^*$} \\
Anaphoric & Pronoun & Alias table & Citation verif. & 1 pass \\
Abbreviated & Partial name & Alias table & Union valid. & 0+2--3\rlap{$^\dagger$} \\
\bottomrule
\end{tabular}
\\[2pt]
{\scriptsize $^*$2 extraction + 2--3 validation (generic-mention detection for ambiguous names, then 2 voting passes).
$^\dagger$0 extraction (deterministic) + 2--3 validation.
$^\ddagger$Extraction needs no knowledge (runs parallel with Layer~1);
validation uses the alias table once Layer~1 completes.
Stratified valid.: $\cup$ for alias-confirmed, $\cap$ for exact-name matches.
Union valid.: all candidates are alias-confirmed, so $\cup$ only.}
\end{table}

% ---------------------------------------------------------------------------
% ============================================================================
% TikZ figure for the pipeline DAG
% Requires: \usepackage{tikz}, \usetikzlibrary{positioning,arrows.meta,fit,backgrounds}
% ============================================================================

\begin{figure*}[t]
\centering
\begin{tikzpicture}[
    node distance=0.8cm and 0.6cm,
    phase/.style={draw, rounded corners, minimum width=2.2cm, minimum height=0.7cm,
                  font=\footnotesize, align=center},
    layer1/.style={phase, fill=blue!8},
    layer2/.style={phase, fill=orange!10},
    layer3/.style={phase, fill=green!8},
    arr/.style={-{Stealth[length=2mm]}, thick, gray!70},
    layerbox/.style={draw=gray!40, dashed, rounded corners=3pt, inner sep=6pt},
    layerlabel/.style={font=\scriptsize\sffamily, gray!60},
]

% Layer 1a: Independent analyses (top row)
\node[layer1] (model)   {Model\\Analysis};
\node[layer1, right=1.2cm of model] (doc) {Document\\Analysis};

% Layer 1b: Enrichment (needs 1a results)
\node[layer1, below=0.7cm of model, xshift=1.7cm] (enrich) {Partial-Ref.\\Refinement};

% Layer 1 internal dependencies
\draw[arr] (model.south east) -- (enrich.north west);
\draw[arr] (doc.south west) -- (enrich.north east);

% Layer 2: Link Recovery
\node[layer2, below=1.4cm of enrich, xshift=-2.0cm] (explicit) {Explicit\\Ref.\ Extr.};
\node[layer2, right=0.6cm of explicit] (contextual) {Contextual\\Ref.\ Disc.};
\node[layer2, right=0.6cm of contextual] (anaphoric) {Anaphoric\\Ref.\ Res.};
\node[layer2, right=0.6cm of anaphoric] (abbreviated) {Abbreviated\\Ref.\ Match};

% Knowledge flow arrows (Layer 1 -> Layer 2)
\draw[arr] (enrich.south west) -- (explicit.north east);
\draw[arr] (model.south) -- (contextual.north west);
\draw[arr] (enrich.south) -- (contextual.north);
\draw[arr] (enrich.south east) -- (anaphoric.north);
\draw[arr] (enrich.south east) -- (abbreviated.north west);

% Layer 3: Link Consolidation
\node[layer3, below=1.4cm of contextual, xshift=0.8cm] (merge) {Merge};

% Dedup dependency: abbreviated matching needs prior strategies' results
\draw[arr, dashed] (explicit.east) -- (abbreviated.west);

% Layer 2 -> Layer 3
\draw[arr] (explicit.south) -- (merge.north west);
\draw[arr] (contextual.south) -- (merge.north);
\draw[arr] (anaphoric.south) -- (merge.north);
\draw[arr] (abbreviated.south) -- (merge.north east);

% Output
\node[below=0.4cm of merge, font=\footnotesize\itshape] (out) {Trace Links $L$};
\draw[arr] (merge) -- (out);

% Layer labels
\node[layerlabel, left=0.3cm of model] {Layer 1a};
\node[layerlabel, left=0.3cm of enrich] {Layer 1b};
\node[layerlabel, left=0.3cm of explicit] {Layer 2};
\node[layerlabel, left=0.8cm of merge] {Layer 3};

% Layer boxes
\begin{scope}[on background layer]
    \node[layerbox, fit=(model)(doc)(enrich), label={[layerlabel]above:Knowledge Acquisition}] {};
    \node[layerbox, fit=(explicit)(abbreviated), label={[layerlabel]above left:Link Recovery}] {};
    \node[layerbox, fit=(merge), label={[layerlabel]above left:Link Consolidation}] {};
\end{scope}

\end{tikzpicture}
\caption{Pipeline architecture as a staged DAG. Layer~1 acquires knowledge
in two sub-tiers: model analysis and document analysis run
concurrently~(1a); partial-reference refinement depends on their
results~(1b).
Explicit reference \emph{extraction} has no knowledge dependencies
and starts concurrently with Layer~1 (shown here under Link Recovery
for grouping); its \emph{validation} uses the alias table once
Layer~1 completes.
Layer~2 runs four recovery strategies; abbreviated matching
depends on prior strategies for deduplication (dashed arrow).
Layer~3 deduplicates their output.}
\label{fig:pipeline}
\end{figure*}

% ---------------------------------------------------------------------------

\subsection{Knowledge Acquisition}\label{sec:knowledge}

Layer~1 builds the structured knowledge that all recovery strategies
consume.
Three analyses produce this knowledge
(Figure~\ref{fig:pipeline}):
Model Analysis classifies component name ambiguity
(\S\ref{sec:names}) and Document Analysis discovers aliases
(\S\ref{sec:alias}), both running in parallel;
Partial-Reference Refinement then enriches the alias table with
LLM-classified partial mappings once both complete.

\subsubsection{Name Ambiguity Classification}\label{sec:names}

Component names vary widely in how specific they are.
Multi-word identifiers like \texttt{ImageProvider} clearly name
architectural roles, but single common English words like
\texttt{Cache} or \texttt{Facade} often appear in text without
referring to any component.
The pipeline must identify which names need special care before
encountering them in the document.

Structural cues---CamelCase boundaries, multiple words, all-uppercase
letters---resolve most cases without LLM involvement: names with these
cues are classified as \emph{architectural}.
The pipeline submits remaining single-word names to the LLM with
few-shot examples for \emph{ambiguous} classification.
The result feeds downstream decisions:
which candidates require generic-mention detection
(\S\ref{sec:entity}) and
which trailing words need capitalized matching
(\S\ref{sec:alias}).

\subsubsection{Alternative Name Discovery}\label{sec:alias}

Documents refer to components by trailing word (\emph{Server} for
\texttt{HTML5~Server}), abbreviation, or synonym.
Without an \emph{alias table} mapping these alternative names to
canonical component names, any matching strategy will
systematically miss such references, limiting recall.
However, LLM-generated alias tables are unreliable: the model
may hallucinate mappings or accept generic words as aliases.

Different alias types offer different validation opportunities.
Abbreviations carry structural evidence---their letters match the
component name's initials---enabling rule-based auto-approval.
Synonyms lack such structure and need semantic judgment, but
CamelCase identifiers can still be auto-approved since they are
constructed names unlikely to be ordinary English.
Partial references are the hardest case: the same trailing word may be
a component reference in one context and a generic noun in another,
so validation requires examining the document contexts where the word
appears.

\paragraph{Core discovery.}
An LLM extracts three alias categories: abbreviations
(parenthetical definitions), synonyms (unambiguous alternative
names), and partial references (trailing words of multi-word names
used alone).
A calibrated few-shot judge validates the discovered mappings:
the judge is instructed to auto-approve abbreviations whose letters
match the component name's initials and CamelCase identifiers;
remaining candidates are judged semantically, rejecting generic
terms (\emph{system}, \emph{process}).
The resulting alias table is consumed by all four recovery strategies.

\paragraph{Partial-reference refinement.}
An LLM-based enrichment step classifies each candidate trailing
word as a standalone entity reference (\textsc{name}) or as an
ordinary English word (\textsc{ordinary}) by examining the
sentences where it appears without the full component name.
A candidate is added to the alias table only if no other
component shares that trailing word and the LLM classifies
its usage as \textsc{name};
ambiguous trailing words require capitalized occurrences
to be considered.

% ---------------------------------------------------------------------------

\subsection{Link Recovery}\label{sec:recovery}

With ambiguity classifications and the enriched alias table in hand,
four strategies target distinct reference types on this shared
knowledge base.
The first three run in parallel; abbreviated matching runs after
them to avoid re-discovering already-linked pairs.

\subsubsection{Explicit Reference Extraction}\label{sec:seed}

Canonical-name matching still risks false positives from ambiguous
names and namespace paths, and a single LLM pass produces
inconsistent results across runs.
Exact-name matches are high-confidence and reliably found by a
single pass, while non-exact matches risk hallucination and benefit
from independent agreement.
This motivates \emph{asymmetric voting} rather than uniform
intersection (which would lose reliable matches) or uniform union
(which would keep hallucinations).

A self-contained extractor (which we call \textbf{ILinker2}) uses two
complementary LLM passes over batched document segments:
Pass~A (mention-framed) asks which components are explicitly
mentioned; Pass~B (actor-framed) asks which component is the
primary actor.
Both passes reject names inside hierarchical namespace paths and
generic English words.
Exact-name matches from \emph{either} pass are accepted (union);
non-exact matches require agreement from \emph{both}
(intersection).
Because ILinker2 requires no Layer~1 knowledge, its extraction
runs in parallel with knowledge acquisition.
Once Layer~1 completes, the raw ILinker2 results are
re-validated through the same evidence-stratified voting used by
contextual discovery (\S\ref{sec:entity}), using the alias table
to distinguish alias-confirmed from exact-name matches.
This additional validation filters false positives that
ILinker2's internal voting missed; any true positives removed
are independently recoverable by the contextual or anaphoric
strategies.

\subsubsection{Contextual Reference Discovery}\label{sec:entity}

Alias-based and functional references require knowledge that the
explicit extractor lacks, but broadening the extraction scope also
increases the risk of hallucinated links.
Running the \emph{same} broad-scope prompt twice independently
and keeping only shared results filters out hallucinations while
preserving the recall benefit of broader scope.
Using two different prompts would introduce prompt-specific biases;
using the same prompt ensures that disagreements between the two
runs reflect only LLM non-determinism.

Knowledge-enriched prompts receive the known aliases from the
alias table, enabling recognition of abbreviations, synonyms, and
partial names that the explicit extractor cannot match.
The extraction scope includes space-separated compound forms,
functional descriptions by name or role, interaction participants
(``X sends data to Y'' covers both X and Y),
passive and prepositional mentions (``handled by X'', ``via X'',
``through X''), and known alias occurrences.
The same prompt runs twice independently, and the pipeline keeps
only candidates found in both passes.

\paragraph{Validation.}
Candidates undergo two-stage validation:
\begin{enumerate}
    \item \emph{Generic-mention detection:}
      For components whose names were classified as ambiguous by
      Model Analysis (\S\ref{sec:names}), an LLM determines
      whether each lowercase occurrence is a component reference or
      ordinary English usage, given the surrounding sentence context;
      the pipeline rejects ordinary-usage occurrences.
    \item \emph{Evidence-stratified voting:}
      Two independent LLM passes with complementary
      foci---\emph{actor-role} (is the component the subject?)
      and \emph{direct-reference} (does the sentence reference the
      component at all?)---validate remaining candidates.
      When a candidate was matched via a known alias, the validation
      prompt includes the alias relationship (e.g., ``Server is a
      known partial reference for HTML5~Server''), and a single
      approving pass suffices (union).
      Exact-name matches require agreement from both passes
      (intersection).
      This stratified rule reflects that alias-confirmed
      matches carry prior knowledge from Layer~1 that justifies
      a lower agreement threshold.
\end{enumerate}

\subsubsection{Anaphoric Reference Resolution}\label{sec:coref}

Sentences containing pronouns (\emph{it}, \emph{they}, \emph{this},
\emph{these}, \emph{that}, \emph{those}, \emph{its}, \emph{their})
carry valid trace information but contain no component name, so
the preceding strategies cannot find them.
However, LLMs can fabricate antecedents---claiming a pronoun refers
to a component that was never mentioned in nearby sentences.
If the LLM is required to cite a specific antecedent sentence,
such fabrications become detectable: the cited sentence either
contains the component name (or a known alias) or it does not.
A separate validation LLM call would itself be subject to
fabrication; deterministic string verification eliminates this risk.

Each pronoun-bearing sentence is presented with a $\pm5$-sentence
context window.
% Resolution constraint justified by Lappin & Leass 1994 (86% accuracy
% with current + 4 preceding sentences) and Centering Theory's
% local-coherence model (Grosz et al. 1995). +-5 context window is
% broader than the 3-sentence resolution bound to aid disambiguation.
The broader window provides surrounding text for disambiguation,
while the prompt instructs the LLM to resolve only within the
previous 3~sentences---a conservative bound that reduces
long-distance fabrications.
The LLM must cite a specific antecedent sentence number.
The pipeline verifies this citation by checking that the
component name (or a known alias) appears verbatim in the
cited antecedent sentence.
Because this check is strict and deterministic, anaphoric
links that pass are accepted without further LLM review.

\subsubsection{Abbreviated Reference Matching}\label{sec:partial}

Contextual discovery may miss sentences where a partial name
appears without other clues in the surrounding text.
Partial names are short strings with high surface-level match rates,
making LLM extraction unreliable: it tends to either over-generate
or miss occurrences.
Deterministic scanning guarantees complete recall for known
partials, and evidence-stratified validation from
\S\ref{sec:entity}---which judges concrete sentence--component
pairs rather than performing open-ended extraction---is a simpler
task where LLMs are more reliable.

The pipeline scans all sentences for word-boundary matches of each
partial reference from the alias table, producing candidates for
sentence--component pairs not already found by earlier strategies.
A clean-mention check ensures the partial name appears as a
standalone word---not inside a dotted path, hyphenated compound, or
qualified name.
Since every partial candidate is alias-confirmed by construction,
all receive union voting through the same evidence-stratified
validation used in contextual discovery (\S\ref{sec:entity}).

% ---------------------------------------------------------------------------

\subsection{Link Consolidation}\label{sec:merge}

The four recovery strategies produce independent candidate
sets, each already validated by its own consensus pattern
(Table~\ref{tab:strategies}).
Because strategies run independently, they may discover the
same $(s_i, c_j)$ pair through different evidence.
Overlaps are removed by first-seen priority:
candidates are processed in strategy order
(explicit, contextual, anaphoric, abbreviated),
and the first source to discover a pair is kept.
